% concatenated from BaburinBoundPaper.tex
\documentclass[11pt,a4paper]{article}
\usepackage{amssymb, amstext, amscd, amsmath, color}
\usepackage{graphicx, url, hhline}
\usepackage{mathrsfs} % supports only capital letters
\usepackage[left=2.50cm, right=2.00cm, top=2.30cm, bottom=2.00cm]{geometry}

\usepackage{float}
%\usepackage{upgreek}

\usepackage{longtable}
\usepackage{multicol}

%\usepackage[T1]{fontenc} % needed for \guillemotright 

\parindent 0.5cm

\hyphenation{module}
\hyphenation{crys-tallo-graphy}
\hyphenation{equi-va-len-ce}
\hyphenation{inequi-va-lent}
\hyphenation{ge-ne-ra-tors}
\hyphenation{ge-ne-ra-ting}
\hyphenation{pro-per-ty}
\hyphenation{nano-structures}
\hyphenation{auto-morphisms}
\hyphenation{auto-morphism}
\hyphenation{nontri-vial}
\hyphenation{re-gu-lar}
\hyphenation{Punkt-lagen}

\newtheorem{defin}{}
\newtheorem{saetze}[defin]{}
\newtheorem{conjec}[defin]{}
\newtheorem{lemmas}[defin]{}
\newtheorem{folger}[defin]{}
\newtheorem{bemerk}[defin]{}
\newtheorem{prop}{Proposition} %Igor added
\newtheorem{theorem}{Theorem}
\newtheorem{corollary}{Corollary}

%\newenvironment{theorem}  {\begin{saetze}\it {\bf Theorem:}}{\end{saetze}}
\newenvironment{conjecture}{\begin{conjec}\it {\bf Conjecture:}}{\end{conjec}}
\newenvironment{lemma}    {\begin{lemmas}\it {\bf Lemma:}}{\end{lemmas}}
%\newenvironment{corollary}{\begin{folger}\it {\bf Corollary:}}{\end{folger}}
\newenvironment{remark}   {\begin{bemerk}\rm {\it Remark:}}{\end{bemerk}}
\newenvironment{proof}    {{\noindent \it Proof}:}{{\fillbox \medskip}} %igor changed
\newenvironment{comments} {{\it Comments}:}{{\hfill \fillbox \bigskip}}

\newcommand{\myref}[1]{[1.4(#1)]}
\newcommand{\fillbox}{\mbox{$\bullet$}}
\newcommand{\ra}{\rightarrow}
\newcommand{\Ra}{\Rightarrow}
\newcommand{\ms}{\mapsto}
\newcommand{\ol}{\overline}
\newcommand{\ti}{\tilde}
\newcommand{\N}{\mathbb N}
\newcommand{\F}{\mathbb F}
\newcommand{\Z}{\mathbb Z}
\newcommand{\Q}{\mathbb Q}
\newcommand{\R}{\mathbb R}
\newcommand{\GAP}{\mathsf{GAP}}
\newcommand{\MAGMA}{\mathsf{MAGMA}}
\newcommand{\bolit}{\boldsymbol}

\newcommand{\sgn}{\operatorname{sgn}}
\newcommand{\rank}{\operatorname{rank}}
\newcommand{\Isom}{\operatorname{Isom}}
\newcommand{\qIsom}{\operatorname{q-Isom}}
\newcommand{\sIsom}{\operatorname{s-Isom}}
\newcommand{\grid}{\operatorname{grid}}
\newcommand{\ggk}{\operatorname{ggk}}
\newcommand{\lgk}{\operatorname{lgk}}
\newcommand{\hc}{\operatorname{hc}}
\newcommand{\ecl}{\operatorname{ecl}}
\newcommand{\Aut}{\operatorname{Aut}}
\newcommand{\sep}{\operatorname{sep}}
\newcommand{\ul}{\underline}
\newcommand{\ec}{\operatorname{ec}}
\newcommand{\lec}{\operatorname{lec}}
\newcommand{\Aff}{\operatorname{Aff}}
\newcommand{\stab}{\operatorname{stab}}
\newcommand{\hxl}{\operatorname{hxl}}
\newcommand{\sql}{\operatorname{sql}}
\newcommand{\cyc}{\operatorname{cyc}}
\newcommand{\pen}{\operatorname{pen}}
\newcommand{\cat}{\operatorname{cat}}
\newcommand{\squeezeup}{\vspace{-2.5mm}} %Igor added

\newenvironment{items}{\begin{list}{$\alph{item})$}
{\labelwidth18pt \leftmargin18pt \topsep3pt \itemsep1pt \parsep0pt}}
{\end{list}}
\newcommand{\bulit}{\item[$\bullet$]}

\begin{document}

\bigskip
\bigskip

\centerline{\large \textbf{Subperiodic groups and bounded automorphisms of periodic graphs}}

\bigskip
\bigskip

\centerline{\large Igor A. Baburin}
\centerline{\emph{Email: baburinssu@gmail.com}}

\bigskip

\noindent A subperiodic group is a group of motions of $d$-dimensional Euclidean space $\R^d$ which contains a translation lattice $\Z^r$ of rank $r < d$ as a subgroup of finite index. A classification into abstract group isomorphism classes is performed for subperiodic groups in dimension~3: 75 \emph{crystallographic} rod groups ($r=1$) and 80 layer groups ($r=2$) are shown to belong to 32 and 34 isomorphism classes, respectively. An easy-to-compute set of invariants is developed for recognizing these isomorphism classes from finite presentations which makes use only of the number of subgroups up to a given finite index~$n$ ($n \leq 12$ for rod groups and $n \leq 8$ for layer groups) and how many of them are normal. Cayley graphs of rod and layer groups are used to illustrate the concept of bounded automorphisms of finite order, \emph{i.e.} those when the distance between a graph vertex and its image has an upper bound. It is proven that a Cayley graph of a crystallographic space group $G$ (in which case $r=d$) possesses bounded automorphisms of finite order, if and only if the respective inverse-closed generating set is stabilized by conjugation by an element of finite order in $G$. As an application, subperiodic groups in $\R^4$ with a three-dimensional translation lattice are used to systematically derive embeddings of three-periodic \emph{ladder graphs} in~$\R^3$.

\section{Introduction}

\bigskip

Crystallographic point groups and space groups are the isometry groups most commonly used in crystallography. In algebraic terms a crystallographic \emph{space group} $G$ is an extension of a free abelian group $T \cong \Z^d$ of rank $d$ by a finite group $P = G/T$ so that $P$ acts faithfully on $T$. The group $P$ is the point group of $G$, and $T$ is its translation subgroup (translation lattice) that we shall also denote by $P(G)$ and $T(G)$, respectively. A \emph{subperiodic group} corresponds to the situation when the rank of the translation lattice $T \cong \Z^r$ to be extended by $P$ is smaller than~$d$ \cite{Koe80}. In this case the group $P$ need not act faithfully on $T$ that gives rise to certain properties which do not hold for space groups. Let $H$ be a subperiodic group in $\R^d$ with translation lattice~$\Z^r$. Then

\begin{enumerate}
\setlength\itemsep{-0.02in}
\item[(a)] $H$ can be isomorphic to some space group in $\R^r$,
\item[(b)] $H$ can have non-trivial finite normal subgroups [owing to the non-faithful action of $P(H)$ on $T(H)$] and therefore cannot be isomorphic to a space group of $\R^d$ for any $d$,
\item[(c)] subperiodic groups can be isomorphic without being affinely equivalent (\emph{e.g.} groups with non-isomorphic point groups can be isomorphic),
\end{enumerate}

\noindent where affine equivalence of groups means, as usual, their conjugacy in the group of all affine mappings of~$\R^d$.

\medskip

\noindent \emph{Example}: Consider the following extensions of a 1-dimensional lattice by finite groups generated by mirror reflections in coordinate planes and 2-fold rotations about coordinate axes: $H_1 = p2mm = \langle m_{xy}, m_{xz}, t_z \rangle$, $H_2 = pmm2 = \langle m_{xz}, m_{yz}, t_z \rangle$ and $H_3 = p222 = \langle 2_{x}, 2_{y}, t_z \rangle$, where $t_z$  stands for the unit translation along the $z$-axis. The non-isomorphism of $H_1$ and $H_2$ can be seen if their maximal finite normal subgroups are regarded: such a subgroup of $H_1$ is generated by one reflection whereas that of $H_2$ is generated by two reflections. As a result, the following factorisations are obtained: $H_1 \cong H_3 \cong D_{\infty} \times C_{2}$ and $H_2 \cong C_{\infty} \times D_{2}$.


\medskip


Subperiodic groups in dimension 3 with translation lattices of rank 1 and 2 are catalogued in Volume~E of the \emph{International tables for crystallography} (2006) where conventional coordinate description is given for 75~\emph{crystallographic} rod groups ($r=1$) and 80~layer groups ($r=2$). The numbers of groups originate from the conjugacy in the group of all \emph{orientation-preserving} affine mappings of $\R^3$ -- \emph {i.e.} from the same equivalence relation that yields 230 space groups in $\R^3$. Note that the finite number of rod groups is only due to the restriction on orders of rotation axes (1, 2, 3, 4, 6) that is arbitrary for $r=1$.

Subperiodic groups can be also viewed as quotient groups of space groups in $\R^d$ with respect to (maximal) normal translation subgroups of rank $r < d$. For example, every hexagonal space group~$G$ in $\R^3$ has normal rank-1 and rank-2 translation subgroups generated by unit translations: $\langle t_z \rangle$ and $\langle t_x, t_y \rangle$. The respective quotients naturally give rise to the rod group $G/\langle t_x, t_y \rangle$ and layer group $G/\langle t_z \rangle$. 

Subperiodic groups of $\R^3$ find applications in solid-state physics as symmetry groups of nanostructures (nanotubes and two-dimensional materials) where representation theory is the main tool \cite{Linegr10}. In this paper we shall develop a more abstract view of subperiodic groups motivated by application to automorphisms of vertex-transitive periodic graphs, specifically \emph{bounded} automorphisms of finite order. The history of the problem goes back to the beginning of 1980s when M.~Gromov proved his famous theorem on groups of polynomial growth \cite{Gromov81} (see Appendix~A for more details). Shortly after that, by extending this result, V.~I.~Trofimov \cite{Trofimov84} characterized vertex-transitive graphs of polynomial growth (of which periodic graphs are a special case) in terms of the structure of automorphism groups, and in terms of partitions of vertex sets with (so to speak) convenient properties. The structures usually dealt with in crystallography have translational symmetry that already implies polynomial growth (see below for the definition). Trofimov's theorem \cite{Trofimov84} thus tells us which additional automorphisms -- beyond the space-group symmetry -- periodic structures might possess. It is precisely where bounded automorphisms of finite order come into play. Eon (2005) rediscovered the idea of bounded automorphisms independently and introduced it to crystallography \cite{Eon2005}. Recently (\emph{ca.}~2019--2023), the results of \cite{Trofimov84} have been proved (or reinterpreted) in the language of topological groups \cite{Salle19, Tes21} that allowed to obtain some new, including also more specialized, results.

In physics and chemistry it is often the case that certain mathematical results have to be \emph{felt} in addition to be rigorously \emph{proven}. In our opinion, bounded automorphisms of finite order in periodic graphs still remain somewhat obscure to crystallographers, and subperiodic groups provide an adequate framework for a better understanding of this phenomenon. Broadly speaking, this paper is a further adaptation of \cite{Trofimov84} to periodic graphs, with the main emphasis on computational group-theoretic aspects. A more abstract treatment of subperiodic groups in relation to graph automorphisms also leads to interesting applications, for example, exploring non-isotopic embeddings of graphs in $\R^d$ which are used to complement a group-theoretic discussion. %throughout the paper to complem / exposition. %constructing

The structure of the paper is as follows. Section~2 contains preliminaries on the terminology used. Section 3 deals with the classification of subperiodic groups in $\R^3$ into abstract group isomorphism classes. The language of generating sets and associated Cayley graphs is used throughout, also for illustrative purposes. Cayley graphs are always characterized with respect to bounded automorphisms of finite order. The logic of our exposition follows from that bounded automorphisms of finite order naturally become more and more rare on going from rod groups to space groups. Section~4 discusses in depth the automorphisms of Cayley graphs which are induced by the conjugation action of a group. Finally, for Cayley graphs of space groups a group-theoretic criterion is proved that is necessary and sufficient for deciding whether or not its full automorphism group is isomorphic to a space group. Section~4 includes various examples and describes some applications, mainly inspired from the analysis of sphere-packing graphs. %related to piecewise-linear embeddings of periodic graphs in $\R^3$.


\section{Preliminaries on groups and graphs}

Our terminology for crystallographic groups is mostly standard (\emph{cf.} above). More general things shall be introduced here. For a group $G$, its center is denoted by $Z(G)$. For subgroup $H<G$, $C_G(H)$ and $N_G(H)$ are its centralizer and normalizer in $G$, respectively. The index of $H$ in $G$ is $|G:H|$. For an inverse-closed subset $S \subset G$, define its stabilizer in $G$ as $Stab_G(S) = \{ g|g^{-1}sg \in S, s \in S, g \in G \}$.

A group is called an FC-group if conjugacy classes of all its elements are finite \cite{Tom84}. The~subgroup $H \leq G$ containing all the elements with finitely many conjugates is the FC-center of $G$. For example, the FC-center of a space group is its maximal translation subgroup.

If $\Gamma$ is a graph, it is said to be \emph{locally finite} if valencies of all its vertices are finite. The vertex set of $\Gamma$ is denoted by $V(\Gamma)$. The distance between vertices $x$ and $y$ in a (connected) graph, $d(x, y)$, is the number of edges in a shortest path (a geodesic) connecting them. The neighbourhood of vertex~$x$ is defined as $\{y|d(x, y) = 1\}$. The automorphisms of $\Gamma$ are understood as adjacency-preserving permutations on $V(\Gamma)$. The group of all automorphisms of $\Gamma$ is denoted by $Aut(\Gamma)$. For $G \leq Aut(\Gamma)$ the stabilizer of a vertex $x$ is defined as usual as $Stab_G(x) = \{g|x=xg, g \in G\}$. If $\Gamma$ is a vertex-transitive graph, then a group $G \leq Aut(\Gamma)$ acting \emph{freely} and \emph{transitively} on $V(\Gamma)$ is called a \emph{regular} group of automorphisms. 

Given a group $G$ and its finite inverse-closed generating set $S = S^{-1} \subset G \setminus \{1_G\}$, recall that a \emph{Cayley graph} of $G$ with respect to $S$ is a simple unidrected graph whose vertices are elements of $G$, and vertices $(g, h)$ are connected by an edge if (and only if) $g^{-1}h \in S$. The action of $G$ on itself by left multiplication induces a regular automorphism group of a Cayley graph \cite{Biggs93}.

An automorphism $g \in Aut(\Gamma)$ is termed \emph{bounded} if there is a constant $C < \infty$ (depending on $g$) such that $d(x, xg) < C$ holds for every $x \in V(\Gamma)$. For Cayley graphs of groups (when $V(\Gamma) = G$), this is equivalent to saying that the conjugacy class of $g$ is finite: since $d(x, gx) = d(e, x^{-1}gx) < C$, the latter holds, if and only if $\{x^{-1}gx, x \in G\}$ is finite ($e=xx^{-1}$). Bounded automorphisms form a normal subgroup $B(\Gamma)$ of $Aut(\Gamma)$ \cite{Seif89all, Seif91}.

The graphs which arise in the context of our paper are $r$-periodic graphs. They are simple undirected graphs with a finite number of vertex- and edge-orbits under the free action of $\Z^r$ ($1 \leq r \leq d$) \cite{Delgado2004}. Translations are their bounded automorphisms of infinite order. This in turn implies that vertex neighbourhoods in a graph, $B(x, n) = \{y|d(x, y) \leq n\}$, grow polynomially: $|B(x, n)| \leq c(x)n^r$ for every $n \in \N$ where $c(x)$ is a constant that becomes independent of $x$ for vertex-transitive graphs \cite{Gods80, Seif97}.

When discussing piecewise linear embeddings of $r$-periodic graphs in $\R^3$, we assume that vertices are mapped to points in $\R^3$, and  edges to straight-line segments. The edges intersect at most at common vertices. Two embeddings are called isotopic if (speaking informally) they can be continuously deformed one into another without breaking edges avoiding also edge intersections across the deformation path. In many cases the graphs we consider are contact graphs of sphere packings (whose vertices are spheres and edges correspond to sphere contacts) which we shall call \emph{sphere-packing graphs} as common in crystallography.


\section{Subperiodic groups in $\R^3$ -- isomorphism classes}

\subsection{Rod groups} \label{rod1}

Table~\ref{t:table1} provides a classification of rod groups into 32 isomorphism classes along with some information on the subgroup structure. Each row corresponds to one isomorphism class, groups are designated following \cite{ITE}.  The relatively simple structure of rod groups often permits to prove group isomorphism by hand, especially if the groups to be compared can be factorized as direct products:

\begin{center}
\begin{tabular}{l}
$p\bar3m1 = p\bar1 \times C_{3v}$ \\ 
$p\bar6m2 = p11m \times C_{3v}$ \\
$p6_3/mmc = p112_1/m \times C_{3v}$,
\end{tabular}
\end{center}

\noindent where $p\bar1 \cong p11m \cong p112_1/m \cong D_{\infty}$, or if an isomorphic mapping of group generators can be easily seen. In many cases a convenient choice of generators can be made from the tabulation of subperiodic homogeneous sphere packings by Koch and Fischer (1978). For checking purposes group isomorphism was proven from group presentations as illustrated by the two characteristic examples below.


The most unexpected group isomorphism is found for the pair ($p4_2/m$, $p\bar42c$). Given the generators %observed/found

\begin{center}
\begin{tabular}{l}
$a = \bar4_z(000) : y, -x, -z$,\\
$b = m(x, y, -1/4) : x, y, -1/2-z$,\\
$c = 2(x, 0, 1/4) : x, -y, 1/2-z$,
\end{tabular}
\end{center}

\noindent we have: $p4_2/m = \langle a, b \rangle$; $p\bar42c = \langle a, c \rangle$. Group presentations on these generators are identical (\emph{cf.} Fig.~\ref{f:p42m} below):

\begin{center}
\begin{tabular}{l}
$\langle a, b  | a^4, b^2, (aba)^2 \rangle \cong \langle a, c | a^4, c^2, (aca)^2 \rangle$.
\end{tabular}
\end{center}

\noindent We have checked that `analogous' isomorphism does not hold for groups containing principal axes of higher order [\emph{e.g.} ($p6_3/m, p\bar6c2$), ($p8_4/m, p\bar8c2$) \emph{etc.}].

\begin{center}
\begin{figure}[h]
\centering
\includegraphics{Fig1.png}
\caption{A Cayley graph for the group defined by the presentation $\langle a, b|a^4, b^2, (aba)^2 \rangle$ (arrows distinguish between $a$~and~$a^{-1}$). See the text for the edge labels.}
\label{f:p42m}
\vspace{-2em}
\end{figure}
\end{center}


The pair ($p4_2/m$, $p\bar42c$) represents an example where distinct affine realisations of an abstract group act on the same embedding of the corresponding Cayley graph in $\R^3$ (Fig.~\ref{f:p42m}). The pairs of isomorphic groups ($p222, p4_222$), ($p321$, $p6_322$) correspond to a different situation that we shall discuss in a more general setting. Consider the following symmetry operations given below by the Seitz symbols [where $R_{\alpha}(\varphi)$ is a shorthand for the counterclockwise rotation by $2\pi/n\, (n \in \N)$ about the  axis $\alpha$]:

%listed below/given below

\begin{center}
\begin{tabular}{l}
$a = \{R_z(\varphi)|000\}, b = \{R_x(\pi)|000\}, c = \{R_x(\pi)|001\}$, \\
$d = \{R_x(\pi)R_z(\varphi)|001\}, f = \{R_x(\pi)R_z(\varphi/2)|00\frac{1}{2}\}$.
\end{tabular}
\end{center}

\noindent Presentations for the two families of groups are as follows ($n \geq 2$)\footnote{For simplicity we write here the same symbol $pn21$, be $n$ even or odd.}: 

\begin{center}
\begin{tabular}{l}
$H_1 = pn21 = \langle a, b, c \rangle = \langle a, b, d \rangle$,\\
$H_2 = p(2n)_n22 = \langle a, b, f \rangle$, \\
$H_1 \cong H_2 \cong \langle a, b, s | a^n, b^2, s^2, (as)^2, (bs)^2 \rangle \, (s=c, d, f)$.
\end{tabular}
\end{center}

\noindent Cayley graphs for these generating sets are isomorphic to quotients of the square lattice with respect to translation subgroup $\langle t_{y}^{n} \rangle$ which we shall denote (following \cite{KF78}) by  $\bf 4^4(0, n)$. Although isomorphic, they show up remarkable differences in their embeddings in $\R^3$: with symmetry $\langle a, b, c \rangle$ `infinite vertical edges' of the graph are disjoint, whereas with symmetry $\langle a, b, d \rangle$ and $\langle a, b, f \rangle$ they are intervowen giving rise to the $n$-tuple helices (Fig.~\ref{f:3tubes}). As a consequence, the two embeddings are non-isotopic. This also manisfests itself in the type of a group-subgroup relation between the symmetry of the graph and the stabilizer of an `infinite vertical edge': the latter is a subgroup of index $n$ in both $H_1$ and $H_2$. However, $S_1 = \langle b, c\rangle \cong p211$ is a \emph{translationengleiche} subgroup of $H_1$ while $S_2 = \langle b, d\rangle \cong \langle b, f \rangle \cong pn_122$  is a \emph{klassengleiche} subgroup of $H_2$ with $n$ times translation period in the $z$-direction. Therefore, cosets of $T(S_2)$ in $T(H_2)$ map individual helices onto each other. A closer look at the embeddings of multiple helices with symmetry $p(2n)_n22$ has shown that this symmetry is optimal in the following sense: it yields vertex-transitive embeddings where edges correspond to shortest equal inter-vertex distances (`bonds'), whereas vertices of different helices could be moved sufficiently far apart. The relevance of such embeddings to crystal structure design has been discussed recently by O'Keeffe and Treacy (2021) \cite{OK21}. % considered structure models for crystal design 

\medskip

\noindent \emph{Remark 1}. The same reasoning applies to generating symmetries of some 1-periodic sphere packings from Koch and Fischer (1978) [designated by $\bf 6^3(0,2)$, $\bf 6^3(4, 2)$, $\bf 4^4(0, 2)$, $\bf 4^4(0, 3)$, $\bf 4^4(0, 4)$] that give rise to non-isotopic embeddings of sphere-packing graphs -- the aspect that remained unnoticed in their paper. For example, for the type $\bf 4^4(0,3)$ they tabulated the generating sets (in the above notation) $\{a, b, c\}$ and $\{a, b, f\}$ for rod groups $p321$ and $p6_322$, resp. Our numerical checks have shown that the set $\{a, b, d\}$ cannot produce a sphere packing of contact number~4. At first sight this discussion might seem as a pure exercise in geometry and group theory. Nevertheless, it appears to be relevant for certain conformations of nanotubes characterised by \emph{intrinsic twist} \cite{Twist2012}. The latter implies that energetically favourable configuration is in general different from that obtained by the geometrical rolling procedure \cite{Evar09} that excludes structures like those of symmetry $p8_422$ and $p6_222$ in Fig.~\ref{f:3tubes}. %Unfortunately, it is to be expected that structures analogous to those in Fig. xx are likely hard to distinguish experimentally.

\begin{center}
\begin{figure}[ht]
\centering
\includegraphics{Fig2.png}
\caption{Embeddings of the $\bf 4^4(0, 4)$ graph with different symmetry (side and top views).}
\label{f:3tubes}
\vspace{-2em}
\end{figure}
\end{center}


\noindent \emph{Remark 2}. Let us explain how the hexagonal group  $p6_222$ may appear as the symmetry group of the `intrinsically tetragonal' graph $\bf 4^4(0,4)$. Note that $p4_222$ is one of its regular groups. Owing to the isomorphism with $p4_222$, the group $p6_222$ is also plausible (Table~\ref{t:table1}). Consider the following group presentation:

\begin{center}
\begin{tabular}{l}
$G = \langle a, b, c, d | a^2, b^2, c^2, d^2, (cd)^2, dacb \rangle$; \\
$p4_222 = \langle 2(xx\frac{1}{4}), 2(xx\frac{\bar1}{4}), 2(x00), 2(0y0)\rangle$; \\
$p6_222 = \langle 2(x00), 2(x\bar{x}\frac{1}{3}), 2(0y\frac{1}{6}), 2(2x, x, \frac{1}{6})\rangle$. 
\end{tabular}
\end{center}

\noindent Here again, $H=\langle a, b\rangle$ is the \emph{translationengleiche} subgroup of index~4 for $p4_222$ ($H \cong p121$ and therefore it corresponds to disjoint `vertical edges'), whereas that for $p6_222$, $H \cong p6_522$, is the \emph{klassengleiche} subgroup of index~4. Hence, the embeddings with symmetry $p8_422$ and $p6_222$ are clearly isotopic (Fig.~\ref{f:3tubes}), so are those of symmetry $p422$ and $p4_222$.


{\tiny
\begin{table}[H]
\begin{center}
\caption{Isomorphism classes of rod groups}
\label{t:table1}
\begin{tabular}{|l|c|l|}
\hline
Groups & Direct product & FC-center \\
\hline
$p1$, $pc11$, $p112_1$, $p4_{1,3}$, $p3_{1,2}$, $p6_{1,5}$ & $C_{\infty}$ & $C_{\infty}$ \\
\hline
$p\bar1, p211, p2/c11, p11m, p112_1/m$  & $D_{\infty}$ & $C_{\infty}$ \\
$p222_1, p2cm, p4_{1,3}22, p3_{1,2}21, p6_{1,5}22$ & & \\
\hline
$pm11$, $p112$, $pcc2$, $pmc2_1$, $p4_2$, $p6_{2,4}$ & $C_{\infty} \times C_2$ & $C_{\infty} \times C_2$ \\
\hline
$p2/m11, p112/m, p222, p2mm$ & $D_{\infty} \times C_2$ & $C_{\infty} \times C_2$ \\
$pccm, pmcm, p4_222, p6_{2,4}22$ & & \\
\hline
$p3, p6_3$ & $C_{\infty} \times C_3$ & $C_{\infty} \times C_3$\\
\hline
$p4$ & $C_{\infty} \times C_4$ & $C_{\infty} \times C_4$ \\
\hline
$p6$ & $C_{\infty} \times C_6$ & $C_{\infty} \times C_6$\\
\hline
$p3m1$, $p6_3mc$ & $C_{\infty} \times D_3$ & $C_{\infty} \times D_3$ \\
\hline
$p\bar3$, $p\bar6$, $p6_3/m$ & $D_{\infty} \times C_3$ & $C_{\infty} \times C_3$\\
\hline
$p4/m$ & $D_{\infty} \times C_4$ & $C_{\infty} \times C_4$ \\
\hline
$p6/m$ & $D_{\infty} \times C_6$ & $C_{\infty} \times C_6$ \\
\hline
$p\bar3m1$, $p\bar6m2$, $p6_3/mmc$ & $D_{\infty} \times D_3$ & $C_{\infty} \times D_3$ \\
\hline
$pmm2$ & $C_{\infty} \times D_2$ & $C_{\infty} \times D_2$ \\
\hline
$p4mm$ & $C_{\infty} \times D_4$ & $C_{\infty} \times D_4$ \\
\hline
$p6mm$ & $C_{\infty} \times D_6$ & $C_{\infty} \times D_6$ \\
\hline
$pmmm$ & $D_{\infty} \times D_2$ & $C_{\infty} \times D_2$ \\
\hline
$p4/mmm$ & $D_{\infty} \times D_4$ & $C_{\infty} \times D_4$ \\
\hline
$p6/mmm$ & $D_{\infty} \times D_6$ & $C_{\infty} \times D_6$\\
\hline
$p\bar4$ &  -- & $C_{\infty} \times C_2$ \\
\hline
$p4_2/m$, $p\bar42c$ & -- & $C_{\infty} \times C_2$ \\
\hline
$p\bar42m$ & -- & $C_{\infty} \times D_2$ \\
\hline
$p3c1$ & -- & $p3c1$ \\
\hline
$p4cc$ & -- & $p4cc$ \\
\hline
$p6cc$ & -- & $p6cc$ \\
\hline
$p4_2mc$ & -- & $p4_2mc$ \\
\hline
$p4_2/mmc$ & -- & $p4_2mc$ \\
\hline
$p\bar3c1$, $p\bar6c2$ & -- & $p3c1$ \\
\hline
$p321$, $p6_322$ & -- & $C_{\infty} \times C_3$ \\
\hline
$p422$ & -- & $C_{\infty} \times C_4$ \\
\hline
$p622$ & -- & $C_{\infty} \times C_6$ \\
\hline
$p4/mcc$ & -- & $p4cc$ \\
\hline
$p6/mcc$ & -- & $p6cc$ \\
\hline
\end{tabular}
\end{center}
\end{table}
}

\subsubsection{Rod groups which are FC-groups}

In Table~\ref{t:table1} rod groups are characterized by FC-centers. For a Cayley graph $\Gamma$ of a group~$G$ we have: $F(G) = G \cap B(\Gamma)$ where $F(G)$ is the FC-center of $G$ and $B(\Gamma)$ is the group of all bounded automorphisms. Some rod groups are their own FC-centers. Apart from abelian groups, infinite FC-groups are very rare in crystallography. A group $G$ is an FC-group if its \emph{commutator subgroup} $G'$ is finite: in this case every element $g \in G$ has indeed finitely many conjugates: $\{h^{-1}gh = g[g, h], h \in G\}$ where $[g, h]=g^{-1}h^{-1}gh \in G'$ \cite{Tom84}. From a crystallographic point of view it is not surprising that all such rod groups are polar, and for them the finiteness of conjugacy classes can be also confirmed from symmetry diagrams \cite{ITE}. 

Cayley graphs of FC-groups provide examples of graphs with \emph{vertex-transitive groups of bounded automorphisms} \cite{Trofimov83, Seif91}, and from this point of view they are generalizations of translation lattices. In this case bounded automorphisms of finite order are realised by rotation(s) about the principal axis and by vertical mirrors, the latter act (informally speaking) `in the directions of missing periodicity'. Bounded automorphisms of infinite order correspond to translations, screw rotations and glide reflections (that are necessarily vertical in rod groups). Moreover, Cayley graphs of FC-groups on one hand, and of $\Z^n$ on the other, have the following characteristic property in common. First observe that in Cayley graphs of translation lattices (or abelian groups in general) the set of elements $\{g | d(x, xg) = 1, g \in G \}$ is the same for any vertex~$x$. In a more general case of FC-groups, a given generating set $S$ can be always extended in such a way that a larger, though finite, generating set $\bar{S}$ is obtained that is invariant under conjugation by any element of the group: $\bar{S} = \{ g^{-1}sg| s \in S, g \in G \}$. In other words, the set $\{s | d(x, xs) = 1, s \in \bar{S} \}$ becomes independent of vertex $x$. %so constructed

\medskip

\noindent \emph{Example:} Consider the FC-group $G = p4cc = \langle s_1, s_2 \rangle$ where $s_1 = \{m_{xxz}|00\frac{1}{2}\}$ and $s_2 = \{m_{yz}|00\frac{1}{2}\}$. The Cayley graph for this generating set is a square lattice rolled along the vector $(4, 4)$. The conjugacy classes of elements obviously correspond one-to-one to those of $G/T(G) \cong D_4$. We have: $S = \{ s_1, s_2\}; \bar{S} = \tilde{S} \cup \tilde{S}^{-1}$, where $\tilde{S} = \{s_1, s_2, s_2s_1s_2, s_1s_2s_1 \}$. The Cayley graph for $\bar{S}$ has valency 8, and therefore does not correspond to a quotient of any planar vertex-transitive graph.

\begin{center}
\begin{figure}[H]
\centering
\includegraphics{Fig3.png}
\caption{Square lattice rolled along the vector $(4, 4)$ (top and side view). The top view is superposed with the symmetry elements of the rod group $p4cc$.}
\label{f:sql44}
\end{figure}
\vspace{-2em}
\end{center}


\subsection{Layer groups}

In Table~\ref{t:table2} layer groups are classified into 34 isomorphism types. This classification parallels that of two-dimensional space groups (\emph{wallpaper groups}). In group-theoretic terms, a layer group is either isomorphic to a wallpaper group, or is a direct product of a wallpaper group and the cyclic group of order~2 (hence the total number of types $34=17+17$). In the first case a layer group is uniquely related to the group-subgroup pair $G/H$ of wallpaper groups with index~2. More specifically, such a layer group which is also referred to as the symmetry group of a two-sided plane, can be constructed by lifting the symmetry operators of $G$ to the $xy$-plane in $\R^3$, and choosing a coset representative of $H$ in $G$ so as it maps opposite sides of the plane onto each other. Group isomorphism between such a layer group and the original wallpaper group $G$ can be verified directly from group presentations as illustrated by Baburin (2026) and Souvignier (2026) \cite{Souv26} in connection with regular groups of the $6^3$-tessellation of the Euclidean plane. The other seventeen layer groups are central extensions of wallpaper groups by the cyclic group of order~2 that is generated by a horisontal mirror reflection. %Since the latter is a central element, it induces the bounded automorphism of finite order in respective vertex-transitive graphs \cite{Gromov83}.

Being in principle simple, the classifiation in Table~\ref{t:table2} seems not to have appeared in the literature before, or anyway the aspect of (abstract) group isomorphism was not emphasized. Some remarks on group isomorphism for subperiodic groups are contained in the papers by Fischer and Koch \cite{KF75, KF78b} who did not go, however, for a full classification beyond a few examples. Very recently, layer groups have been described in terms of orbifolds but again the aspect of group isomorphism was not considered \cite{layer25}.

{\tiny
\begin{table}[t]
\begin{center}
\caption{Isomorphism classes of layer groups}
\label{t:table2}
\begin{tabular}{|l|l|l|l|}
\hline
Type & Groups & Type & Groups\\
\hline
$p1$ & $p1, p11b$ & $p1 \times C_2$ & $p11m$ \\
$p2$ & $p112, p\bar1, p112/b$ & $p2 \times C_2$ & $p112/m$ \\
$pm$ & $p1m1, p121, p2_1ma, p2aa, p2mb, cm2a$ & $pm \times C_2$ & $pm2m$ \\
$pg$ & $p1a1, p12_11, p2_1ab$ & $pg \times C_2$ & $p2_1am$ \\
$cm$ & $c1m1, c121, p2_1mn, p2an$ & $cm \times C_2$ & $cm2m$ \\
$p2mm$ & $pmm2, p12/m1, p222, pmaa, pmma, cmma$ & $p2mm \times C_2$ & $pmmm$ \\
$p2mg$ & $pma2, p12_1/m1, p12/a1, p2_122, pbaa, pmab$ & $p2mg \times C_2$ & $pmam$ \\
$p2gg$ & $pba2, p12_1/a1, p2_12_12$ & $p2gg \times C_2$ & $pbam$ \\
$c2mm$ & $cmm2, c2/m, c222, pban, pbmn, pmmn$ & $c2mm \times C_2$ & $cmmm$ \\
$p4$ & $p4, p\bar4, p4/n$ & $p4 \times C_2$ & $p4/m$ \\
$p4mm$ & $p4mm, p422, p\bar42m, p\bar4m2, p4/nbm, p4/nmm$ & $p4mm \times C_2$ & $p4/mmm$ \\
$p4gm$ & $p4bm, p42_12, p\bar42_1m, p\bar4b2$ & $p4gm \times C_2$ & $p4/mbm$ \\
$p3$ & $p3$ & $p3 \times C_2$ & $p\bar6$ \\
$p3m1$ & $p3m1, p312$ & $p3m1 \times C_2$ & $p\bar6m2$ \\
$p31m$ & $p31m, p321$ & $p31m \times C_2$ & $p\bar62m$ \\
$p6$ & $p6, p\bar3$ & $p6 \times C_2$ & $p6/m$ \\
$p6mm$ & $p6mm, p\bar3m1, p\bar31m, p622$ & $p6mm \times C_2$ & $p6/mmm$ \\
\hline
\end{tabular}
\end{center}
\vspace{-1em}
\end{table}
}

The abelian groups in Table~\ref{t:table2} are $p1, p11b$ and $p11m$ which are also the only FC-groups among layer groups. Since the horizontal mirror is a central element, it induces the bounded automorphism of order 2 in respective vertex-transitive graphs \cite{Gromov83}. The isomorphism between $p1$ and $p11b$ is noteworthy from a crystallographic point view.

As an application, let us point out the relevance of layer groups to non-isotopic embeddings of 2-periodic non-planar graphs in $\R^3$ (\emph{cf.} also \cite{Bab26}). Consider the {\bf{KIa}} graph \cite{KF78} that often arises in crystallography in diffent contexts. Its automorphism group, $Aut(\Gamma)$, is isomorphic to wallpaper group $p4mm$. By group isomorphism, and by considering also the vertex stabilizer, it can be shown that there are exactly two layer groups for its maximum-symmetry embeddings: $p\bar4m2$ and $p4/nbm$ (Fig.~\ref{f:klee}). The {\bf{KIa}} graph is built from two sets of chains which lie separately in parallel horisontal planes in the first case, and are interwoven in the other. Interestingly, the first ever occurrence of this graph in crystal structures (CSD code XAKCIV, sp.gr. $Pnna$) corresponds to the interwoven variant with symmetry $p2an$ (that is an index-4 subgroup of $p4/nbm$) (see Fig.~51 in \cite{Proser03}).

\begin{center}
%\vspace{-0.5em}
\begin{figure}[h]
\centering
\includegraphics{Fig4.png}
\caption{Embeddings of the {\bf{KIa}} graph with different layer-group symmetry.}
\label{f:klee}
\end{figure}
\vspace{-1em}
\end{center}

\subsection{Invariants for subperiodic groups}

Following the ideas of the \textsc{Small Groups} library \cite{Besche2002} (as implemented in $\GAP$ \cite{GAP}), a practical approach for classifying groups consists in developing a set of invariants that is unique for a given isomorphism class, and a case-by-case consideration of groups sharing the same (finite sets of) invariants. It has to be emphasized that such invariants allow to distinguish groups within a given finite set. For rod groups and layer groups, the invariants can be chosen particularly simple. Every group is characterized by a `fingerprint' that is an ordered list of 2-tuples [($m_1,n_1$), ... ($m_k, n_k$)]. The $kth$ tuple ($m_k, n_k$) refers to the number $m_k$ of conjugacy classes of subgroups of index $k+1$, and the number $n_k$ of normal subgroups of the same index. Experiment suggests that the maximal index to be considered is 12 for rod groups and 8 for layer groups, respectively. We emphasize that these invariants (deposited in the Supporting information) can be efficiently obtained from finite presentations using standard tools of computational group theory. The algorithms for identifying wallpaper groups proposed in the literature are usually oriented towards matrix groups and involve a more elaborate case-by-case analysis of the group structure (as done in \cite{Jones11}) or an explicit computation of a conjugating matrix given the `reference' list of groups \cite{Cryst2019}. % in some standard form

The importance of considering normal subgroups can be clarified by the following example. Wallpaper groups $cm$ and $p2gg$ have the same number of conjugacy classes of subgroups up to index~11. Taking into account the number of normal subgroups allows to distiguish them starting from index-3 subgroups:

\begin{center}
\begin{tabular}{l}
$cm: \, \, \, \, [ (3, 3), (2, 1), (7, 3), (2, 1), (7, 4), (2, 1), (13, 5), ... ]$;\\
$p2gg:  [ (3, 3), (2, 0), (7, 3), (2, 0), (7, 2), (2, 0), (13, 3), ... ]$.
\end{tabular}
\end{center}


\section{Automorphisms of Cayley graphs from group actions} % ? auts in graphs

\subsection{Theory} \label{auts}

We begin with the fundamental result on Cayley graph automorphisms induced by the conjugation action of a group.

\begin{prop}
\label{prop1}
\emph{(Biggs, 1993, ch.~16)}. Let $\Gamma$ be a Cayley graph of a group $G$ corresponding to the generating set $S=S^{-1}$. The element $h \in Aut(G)$ such that $h^{-1}Sh = S$ (i.e., $h$ stabilizes $S$ setwise by conjugation) induces an automorphism of $\Gamma$ fixing the vertex representing the identity of $G$ (centralizing elements give rise to trivial automorphisms).
\end{prop}

Let us emphasize that already the action of $G$ on itself by conjugation allows us to find additional automorphisms of $\Gamma$ which are not left shifts of $G$. This is because $S$ can be stabilized by a non-trivial subgroup of $G$, although $G$ is regular on vertices of $\Gamma$ (Table~\ref{t:table3}). Everywhere in the following the stabilizer of $S$ refers to the inversed-closed generating set ($S \cup S^{-1}$).


{\tiny
\begin{table}[t]
%\begin{center}
\caption{Stabilizers of generating sets for some Cayley graphs}
\label{t:table3}
\vspace{0.3em}
\small \raggedright \hspace{3.2em} `--' in the last column refers to trivial $Z(G)$.\\
\vspace{-1em}
\begin{center}
\begin{tabular}{|l|l|l|l|l|}
\hline
Graph & Gen. set $S$ & $G$ & $Stab_G(S)$ & $Stab \in Z(G)$ \\
\hline
$K_{3,3}$  &  $(1, 2), (1, 3), (2, 3)$ & $D_3$ & $D_3$ & -- \\
\hline
\rule{0pt}{2.5ex} trunc. octahedron & $m_{xxz}, m_{x\bar{x}z}, m_{xyy}$ & $\bar43m$ & $C_2 = \langle 2(x00)\rangle$ & -- \\
%\hline
& $4(00z), 2(x0x)$ & $432$ & $C_2 = \langle 2(x00)\rangle$ & -- \\
\hline
\rule{0pt}{2.5ex} rhombocuboctahedron & $3(xxx), m_{xz}, m_{yz}$ & $m\bar3$ & $C_i = \langle \bar1(000) \rangle$ & yes \\
%\hline
& $4(00z), 3(xxx)$ & $432$ & $C_2 = \langle 2(x\bar{x}0)\rangle$ & -- \\
\hline
\rule{0pt}{2.5ex} trunc. tetrahedron & $3(xxx), 2(00z)$ & $23$ & $C_1$ & -- \\
\hline
\rule{0pt}{2.5ex} icosahedron & $3(xxx), 3(xx\bar{x}), 2(x00)$ & $23$ & $C_2 = \langle 2(00z)\rangle$ & -- \\
\hline
\rule{0pt}{2.5ex} $\bf 4^4(0, 3)$ & $3(00z), 2(x00), 2(x0\frac{1}{2})$ & $p321$ & $C_1$ & -- \\
\hline
\rule{0pt}{2.5ex} $\bf 4^4(0, 4)$ & $\bar4_z(00\frac{1}{4}), 2(x00), 2(x0\frac{1}{2})$ & $p\bar42c$ & $C_2 = \langle 2(00z)\rangle$ & yes \\
\hline
\rule{0pt}{2.5ex} $3/10/c1$ ({\bf srs}) & $4_3(\frac{1}{4}\frac{1}{4}z), 2(00z)$ & $I4_1$ & $p1 = \langle t_z \rangle$ & yes \\
\hline
\rule{0pt}{2.5ex} $4/6/c1$ ({\bf dia}) & $\bar1(000), 2(0y\frac{1}{4}), b(\frac{1}{4}yz)$  & $Pbcn$ & $C_1$ & -- \\
\hline
\rule{0pt}{2.5ex} $4/6/t4$ ({\bf cds}) & $4_3(\frac{1}{4}\frac{1}{4}z), b(\frac{1}{4}yz)$  & $I4_1cd$ & $p2_1 = \langle 2_1(\frac{1}{4}\frac{1}{4}z) \rangle$ & no \\
\hline
\rule{0pt}{2.5ex} $6/4/c1$ ({\bf pcu}) & $4_1(00z), t_x, t_y$  & $P4_1$ & $p4_1 = \langle 4_1(00z) \rangle$ & no \\
%\hline
\rule{0pt}{2.5ex} &  $4(00z), 4(\frac{1}{2}\frac{1}{2}z), c(xxz)$ & $P4cc$ & $pc = \langle c(xxz) \rangle$ & no \\
\hline

%\rule{0pt}{2.5ex} {\bf cds} & $4_1(00z), t_x$  & $P4_1$ & $p2_1 = \langle 2_1(00z) \rangle$ & no \\
%\hline
\end{tabular}
\end{center}
\end{table}
}

\medskip

\noindent \emph{Examples}: (a) Following Biggs (1993), construct a Cayley graph $\Gamma$ for $D_3$ given the generating set of three involutions $\{ (1, 2), (1, 3), (2, 3) \}$. Let us determine additional automorphisms of $\Gamma$ (that is isomorphic to $K_{3,3}$ -- see Fig.~\ref{f:K33}) from the conjugation action of $D_3$ on itself (below $v_i$ are vertices of $\Gamma$):

\begin{center}
\begin{tabular}{l}
$G = \{ (1), (1,2), (1,2,3), (2,3), (1,3,2), (1,3) \} \mapsto \{ v_1, v_2, v_3, v_4, v_5, v_6 \}$;\\
$G^{(1, 2)} = \{ (1), (1,2), (1,3,2), (1,3), (1,2,3), (2,3) \} \mapsto \{ v_1, v_2, v_5, v_6, v_3, v_4 \} \cong (3, 5)(4, 6)$;\\
$G^{(1, 3)} = \{ (1), (2,3), (1,3,2), (1,2), (1,2,3), (1,3) \} \mapsto \{ v_1, v_4, v_5, v_2, v_3, v_6 \} \cong (2, 4)(3, 5)$.\\
\end{tabular}
\end{center}

\noindent As result, we have found a larger automorphism group $\ti{G}(\Gamma) \cong D_3 \times D_3$ of order 36 (note that $|Aut(\Gamma)|=72$):  

\begin{center}
\begin{tabular}{l}
$\ti{G}(\Gamma) = \langle (1,2)(3,4)(5,6), (1,4)(2,5)(3,6), (1,6)(2,3)(4,5), (3, 5)(4, 6), (2, 4)(3, 5) \rangle$.
\end{tabular}
\end{center}

\begin{center}
\begin{figure}[H]
\centering
\includegraphics{Fig5.png}
\caption{$K_{3,3}$ as a Cayley graph for $D_3$.}
\label{f:K33}
\vspace{-2em}
\end{figure}
\end{center}

\noindent (b) Generating sets of polar space groups in Table~\ref{t:table3} are characterised by infinite stabilizers which are rod groups (in this case $Stab_G(S) \cap Z(G) \neq \{e\}$). Consider the $6/4/c1$ ({\bf pcu}) graph that is isomorphic to the three-dimensional cubic lattice. The fourfold screw rotation $f$ induces the following permutation of generators:

\begin{center}
\begin{tabular}{l}
$ \{ f, f^{-1}, t_x, t_{x}^{-1}, t_y, t_y^{-1} \} \mapsto \{ f, f^{-1},  t_y, t_y^{-1} , t_{x}^{-1}, t_x \}$;\\
$\hspace{3.65em} \{1, 2, 3, 4, 5, 6\} \mapsto \{1, 2, 5, 6, 4, 3\} \cong (3,5,4,6)$.\\
\end{tabular}
\end{center}

\noindent This permutation induces in turn an automorphism of the graph that is realised by the fourfold rotation $4(00z)$ that leads to the enhancement of the automorphism group from $P4_1$ to $P4$.

\medskip

For an arbitrary group $G$ there can be different `mechanisms' for the additional stabilization of generating sets (\emph{e.g.} \cite{Watk71}). The following proposition describes the situation that we shall encounter most often. %most relevant for crystallographic space groups.  %frequent add.stabilisation to occur

\begin{prop}
\label{prop2}
Let $\Gamma$ be a vertex-transitive locally finite graph, and let $G < Aut(\Gamma)$ be a regular automorphism group of $\Gamma$. Assume that $C_{Aut(\Gamma)}(G)$ is nontrivial. Put $z \in C_{Aut(\Gamma)}(G)$. Then the generating set $S = \{g | d(x, xg) = 1, x \in V(\Gamma), g \in G \}$ of $G$ is stabilized by some $hz \in G$, where $h \in Stab_{Aut(\Gamma)}(x)$. %, h \neq 1$.
\end{prop}

\begin{proof} Only a sketch of the proof is provided here.\\
\noindent 1) For the generating set $S$ of a regular group $G$ we have: $S = h^{-1}Sh$ for some $h \in Stab_{Aut(\Gamma)}(x)$. The element $hz$ stabilizes $S$ in the same way: $(hz)^{-1}S(hz) = h^{-1}Sh$. For a nontrivial stabilization of $S$, $h \neq 1$ and $z \notin G$ must hold. Whether such $hz$ is actually in $G$ cannot be guaranteed in general but may well be the case under more restrictive assumptions as introduced in the following. \\
2) Let $z \in Z(Aut(\Gamma)) \neq \{e\}$. Assume that some $h$ induces a group automorphism of $G$, and $Stab_{Aut(\Gamma)}(x)$ is not self-normalizing in $Aut(\Gamma)$. Then $hz \in N_{Aut(\Gamma)}(Stab_{Aut(\Gamma)}(x))$. The `normalizer of the stabilizer' permutes the vertices fixed by the stabilizer \cite{Wie64}. Since $G$ is transitive, there does exist such $hz \in G$ which maps fixed vertices of the stabilizer onto each other.
\end{proof}

\medskip

\noindent \emph{Example}. Table~\ref{t:table3} contains two instances (rhombocuboctahedron with symmetry $m\bar3$ and $\bf 4^4(0, 4)$) for which $Z(G) = Z(Aut(\Gamma))$ holds, and therefore a non-trivial stabilization of $S$ does not occur.

\bigskip

\noindent If $G$ is a space group, more can be said about the stabilizer of a generating set.

\begin{prop}
\label{prop3}
Let $G$ be a space group and $S = S^{-1} \subset G$ be its finite set of generators. Then the group $H = Stab_G(S)$ is an FC-group.
\end{prop}


\begin{proof} By assumption, $S$ is finite, and therefore the number of elements $h \in H$ \emph{modulo} $Z(H)$ which stabilize $S$ nontrivially is also finite. Hence $|H:Z(H)|$ must be finite that implies $H$ to be an FC-group \cite{Tom84}.
\end{proof}


\begin{prop}
\label{prop4}
Let $\Gamma$ be a Cayley graph of a space group $G$ with respect to the generating set $S = S^{-1}$, and let $Aut(\Gamma)$ be isomorphic to a space group. Then $S$ is stabilized by some $hz \in G$, where $h \in Stab_{Aut(\Gamma)}(x)$ and $z \in C_{Aut(\Gamma)}(G)$. Moreover, $hz$ is either the identity or an element of infinite order.
\end{prop}

\begin{proof} Since $G$ is a subgroup of $Aut(\Gamma)$ of finite index, $z \in C_{Aut(\Gamma)}(G)$ is either the identity or a nontrivial translation. As in Proposition~\ref{prop2}, $S = h^{-1}Sh$ for some $h \in Stab_{Aut(\Gamma)}(x)$. The element $hz$ stabilizes $S$ in the same way $(hz)^{-1}S(hz) = h^{-1}Sh$, if and only if $hz = zh$. If $C_{Aut(\Gamma)}(G) = \{e\}$, then $hz = 1$  since $h = 1$ must hold ($G$ is a regular group of automorphisms). So assume $z \neq 1$. A nontrivial stabilization of $S$ is achieved whenever $h \neq 1$. Then, $hz=zh$ is only possible if $z$ belongs to the invariant subspace of $h$, in which case $hz$ is of infinite order.
\end{proof}

%\medskip

\begin{prop}
\label{prop5}
Under the assumptions of Proposition~\ref{prop4}, the group $Stab_G(S)$ is torsion-free abelian. % $H \leq G$
\end{prop}

\begin{proof} By Proposition~\ref{prop2}, $Stab_G(S)$ is an FC-group that is a subgroup of a space group. Further, by Proposition~\ref{prop4}, it has no nontrivial finite-order elements. Then it is torsion-free abelian.
\end{proof}

\medskip

\noindent Now we are in a position to state our main result.

\begin{theorem}
\label{teo1}
Let $\Gamma$ be a Cayley graph of a space group $G$ with respect to the generating set $S = S^{-1}$. Then the following assertions are equivalent:\\
1) there exists element $g \in G, g \neq 1$ of finite order such that $S = g^{-1}Sg$;\\
2) $\Gamma$ has nontrivial bounded automorphisms of finite order. 
\end{theorem}

\begin{proof} By Proposition~\ref{prop4}, the existence of a stabilizing element $g$ of finite order implies that $Aut(\Gamma)$ is not isomorphic to a space group. At the same time, $\Gamma$ is a graph with polynomial growth since $\Z^d$ acts on $V(\Gamma)$ with finitely many orbits. Then, from the characterization of graphs with polynomial growth \cite{Trofimov84}, it follows that $\Gamma$ has bounded automorphisms of finite order.
\end{proof}

\begin{corollary}
\label{cor1}
The finite-order elements of $G$ stabilizing $S$ form a finite subgroup $H$ of $G$. $V(\Gamma)$ is partitioned into finite blocks which are orbits of $H$ and its conjugates in $G$.
\end{corollary}

\medskip

\noindent The above results provide a recipe for constructing Cayley graphs of space groups which \emph{a priori} possess bounded automorphisms of finite order.

\begin{corollary}
\label{cor2}
Let $G$ be a space group, and $H$ its nontrivial finite subgroup. Let $S=S^{-1}$ be a finite set of generators for $G$ such that $S$ is a union of finitely many orbits of $H$ on elements of $G$. Then, by construction, $H = Stab_G(S)$ and the respective Cayley graph has bounded automorphisms of finite order.
\end{corollary}

\noindent Moreover, the following general construction due to de la Salle and Tessera (2019) (which can be viewed as `symmetrization' of a generating set by a finite subgroup) leads to Cayley graphs with infinite vertex stabilizers in $Aut(\Gamma)$ (an example will be given in the next section).

\begin{theorem}
\label{teo2}
\emph{\cite[Lemma~6.1]{Salle19}} Let $G$ be a space group with a finite generating set $S=S^{-1}$ and $H$ a nontrivial finite subgroup of $G$. The generating set $\{hsh'|h, h' \in H, s\in S \}$ gives rise to a Cayley graph $\Gamma$ with infinite vertex stabilizers in $Aut(\Gamma)$. %avoid HSH In this case $H \leq Stab_{HSH}(G)$.
\end{theorem}

\noindent However, not all space groups allow for Cayley graphs with bounded automorphisms of finite order.

\begin{corollary}
\label{cor3}
\emph{(cf. \cite[Corollary~6.6]{Salle19})} Let $G$ be a space group without nontrivial finite subgroups (it is a fixed-point-free space group that is often termed a Bieberbach group). Any Cayley graph of $G$ with respect to a finite generating set has no bounded automorphisms of finite order.
\end{corollary}


\subsection{Cayley graphs of crystallographic groups: applications} %space groups

The above theory has the following practical implications. Let $\Gamma$ be a Cayley graph of a space group $G$. The conjugation action of $G$ (or possibly of its affine normalizer) can be applied to compute additional automorphisms of $\Gamma$ that complements and extends available methods. The latter include the method of barycentric placement \cite{Delgado2004} and the method of vertex stabilizer computation \cite{Bab20}. The method of barycentric placement can be applied when $Aut(\Gamma)$ is isomorphic to some space group, and furthermore, if neighbourhoods of vertices of the graph are injectively mapped to their barycentric positions (in terms of O.~Delgado such embedding is \emph{neighbour-unique}). The method \cite{Bab20} puts no restrictions on $Aut(\Gamma)$ and allows to compute vertex stabilizer in $Aut(\Gamma)$ as a permutation group from a sufficiently large finite ball that requires, however, more preprocessing steps. Infinite stabilizers can be recognized from a characterstic growth of stabilizer order as soon as the radius of the ball is iteratively enlarged.

The main problem in computing $Aut(\Gamma)$ consists in handling bounded automorphisms of finite order. It is known that if a periodic graph has bounded automorphisms of finite order, then its barycentric placement shows up vertex-collisions. However, this statement cannot be reversed, and there exist periodic graphs with vertex-collisions but with no bounded automorphisms of finite order (\emph{cf.} examples below). Theorem~\ref{teo1} provides a necessary and sufficient criterion when $Aut(\Gamma)$ is a space group for a very broad class of vertex-transitive periodic graphs which are Cayley graphs of space groups. 

The examples in the next section are aimed at exhibiting certain aspects of Cayley graphs which are beyond geometric intuition. To facilitate the discussion, the following uniform notation is introduced: $\Gamma$ is a Cayley graph of a group $G$ for a generating set $S$, $B(\Gamma) \trianglelefteq Aut(\Gamma)$ is the group of all bounded automorphisms of $\Gamma$, and $B_o(\Gamma) \trianglelefteq Aut(\Gamma)$ is the group of all bounded automorphisms of finite order \cite{Seif89all}. Sometimes we shall also use the notation $G_o = G \cap B_o(\Gamma)$. Recall that $Stab_G(S)$ always refers to $Stab_G(S \cup S^{-1})$.

\subsubsection{Examples}

\noindent \emph{1-Periodic graphs}. The tabulation of homogeneous sphere packings with rod-group symmetry by Koch and Fischer (1978) provides a rich collection of Cayley graphs. Therein, four types [$\bf 3^6(0,2)$, $\bf 3^34^2(0,2)$, $\bf 4^4(2,2)$ and $\bf 6^3(0,2)$] stand out which have uncountable automorphism groups. They are all related to the graph $\bf 3^6(0,2)$ (that is the direct product of the infinite chain and $K_2$) which can be visualized as a band of edge-sharing tetrahedra (Fig.~\ref{f:infstab}). This graph is known in the mathematical literature since 1980s as a classical example of a vertex-transitive graph with infinite vertex stabilizers \cite{Jung84, Seif91}. The reason for this is that vertices within the blocks $\{u_i, v_{i} \}$ (Fig.~\ref{f:infstab}) can be permuted independently. That this graph can be also viewed as a Cayley graph of an infinite dihedral group $D_{\infty}$ (one of the two one-dimensional space groups) does not seem to be well-known. 

\begin{center}
\begin{figure}[h]
\centering
\includegraphics{Fig6.png}
\caption{Embeddings of the $\bf 3^6(0,2)$ graph with rod-group symmetry $p4_122$ (top) and $p6_122$ (bottom).}
\label{f:infstab}
\end{figure}
\vspace{-2em}
\end{center}

\begin{center}
\begin{figure}[h]
\centering
\includegraphics{Fig7.png}
\caption{The $(3, 0)$ carbon nanotube.}
\label{f:zig}
\end{figure}
\vspace{-2em}
\end{center}

\noindent Consider the minimal generating set for $D_{\infty}$:

\begin{center}
\begin{tabular}{l}
$D_{\infty} = \langle a, b | a^2, (ab)^2 \rangle$.\\
\end{tabular}
\end{center}

\noindent By Theorem~\ref{teo2}, it can be symmetrized with respect to $\langle a \rangle$ as follows:

\begin{center}
\begin{tabular}{l}
$D_{\infty} = \langle a, b, ab, ba \rangle$.\\
\end{tabular}
\end{center}

\noindent These abstract generators can be realized, for example, by the following isometries (\emph{cf.} Table~1):

\begin{center}
\begin{tabular}{l}
$p2/c  \, \, =  \langle \bar1(000), c(x0z), 2(0y\frac{1}{4}), 2(0y\frac{\bar1}{4}) \rangle$;\\
$p4_122 = \langle 2(x00), 4_1(00z), 2(xx\frac{1}{8}), 2(x\bar{x}\frac{\bar1}{8}) \rangle$;\\
$p6_122 = \langle 2(x00), 6_1(00z), 2(xx\frac{1}{12}), 2(x\bar{x}\frac{\bar1}{12} \rangle$.
\end{tabular}
\end{center}


\noindent The orbits of $\langle a \rangle $ on $V(\Gamma)$ are two-element blocks (corresponding to shared edges of tetrahedra) which are the orbits of $B_o(\Gamma)$. The generating sets for $p2/c$ and $p4_122$ give rise to sphere-packing graphs and therefore are listed in \cite{KF78}; the one for $p6_122$  does not allow for such a realization (Fig.~\ref{f:infstab}). Moreover, this graph also admits a vertex-transitive group $G < B(\Gamma)$:

\begin{center}
\begin{tabular}{l}
$C_{\infty} \times C_2 \cong pcc2 = \langle c(x0z), c(0yz), 2(00z) \rangle$. % ABELIAN FC-group
\end{tabular}
\end{center}

\noindent Here, since $G = pcc2$ is abelian (\emph{cf.} Theorem~5.14 in \cite{Jung84}), the generating set is trivially stabilized by the whole group. $G_o = \langle 2(00z) \rangle$ generates two-element blocks.

1-Periodic graphs show clearly certain phenomena which are hard to imagine otherwise. For example, \emph{zigzag} carbon nanotubes are graphs where  $Stab_{Aut(\Gamma)}(x) < B_o(\Gamma)$, \emph{i.e.} vertices are stabilized only by bounded automorphisms of finite order (\emph{e.g.} by a mirror plane of any boat-shaped 6-cycle in Fig.~\ref{f:zig}). For the $(3, 0)$ tube we have: 

\begin{center}
\begin{tabular}{l}
$Aut(\Gamma) = p6_3/mmc, B(\Gamma) = p6_3mc, B_o(\Gamma) = C_{3v}$,\\
$Stab_{Aut(\Gamma)}(x) = \langle m(2x,x,0) \rangle$.
\end{tabular}
\end{center}

\noindent Note here that vertex stabilizer is a self-normalizing subgroup of $Aut(\Gamma)$ (\emph{cf.} Proposition~\ref{prop2}).


\medskip

\noindent \emph{2-Periodic graphs}. The simplest examples of 2-periodic graphs with bounded automorphisms of finite order are so-called \emph{ladder graphs} which are obtained by linking directly two copies of the same graph. Construct a ladder graph $\Gamma$ from the two square lattices (Fig.~\ref{f:sql}). In this case the automorphism group, $Aut(\Gamma)$, is a layer group:

\begin{center}
\begin{tabular}{l}
$Aut(\Gamma) = p4/mmm, B(\Gamma) = p11m \cong \Z^2 \times C_2, B_o(\Gamma) \cong C_2$,\\
$Stab_{Aut(\Gamma)}(x) = \langle m_{xz}, m_{xxz} \rangle = C_{4v} \cong D_4$.
\end{tabular}
\end{center}


\noindent Owing to relatively rich vertex stabilizer, this graph has 56 regular automorphism groups (up to conjugacy in $Aut(\Gamma)$), 40 of which are isomorphic to the wallpaper groups. As an example, we provide here generating sets for three regular groups which are all isomorphic to the wallpaper group $p4mm$ (\emph{cf.} Table~\ref{t:table2}):

\begin{center}
\begin{tabular}{l}
%$cmma = \langle m(x0z), m(0yz), m(\frac{1}{2}yz), 2(\frac{1}{4}\frac{1}{4}z), 2(\frac{1}{4}y0) \rangle$,
$p422 = \langle 2(xx0), 4(00z), 4(\frac{1}{2}\frac{1}{2}z) \rangle$;\\
$p\bar4m2 =  \langle 2(xx0), m(x0z), m(0yz), m(x\frac{1}{2}z), m(\frac{1}{2}yz) \rangle$;\\
$p4/nbm = \langle 2(x00), 4(00z), m(x, \frac{1}{2}-x, z), m(\frac{1}{2}+x, x, z) \rangle$.
\end{tabular}
\end{center}

\noindent Clearly each set is stabilized by the respective in-plane 2-fold rotation exchanging the two subgraphs. Again, the orbits of these stabilizers are two-element blocks which are in turn the orbits of $B_o(\Gamma)$.

\begin{center}
\begin{figure}[h]
\centering
\includegraphics{Fig8.png}
\caption{A finite portion of the square ladder.}
\label{f:sql}
\end{figure}
\vspace{-2em}
\end{center}

%\medskip

\noindent \emph{3-Periodic graphs}. Following Corollary~\ref{cor2}, Cayley graphs of space groups in any dimension can be constructed so that $B_o(\Gamma) \neq \{e\}$. Such graphs are rare in crystallography mostly because of their undesirable geometrical properties that is, however, not always the case (see next section). Let us generalize the example of $K_{3,3}$ (\emph{cf.} the beginning of Section~\ref{auts}) to space group $G = P321$:

\begin{center}
\begin{tabular}{l}
$S = \langle 2(x00), 2(0y0), 2(xx0), t_x, t_y, t_xt_y, t_z \rangle$;\\ 
$Stab_G(S) \cong D_3$.
\end{tabular}
\end{center}

\noindent The respective Cayley graph has valency 11. All its vertices collide in the barycentric placement to the points of the hexagonal Bravais lattice. The computation of vertex stabilizer \cite{Bab20} shows that $Stab_{Aut(\Gamma)}(x) \cong (C_2)^3 \times D_3 \times D_3$, and hence the index $|Aut(\Gamma) : G| = 288$.

Another generating set for $P321$ can be constructed by extending that of the $\bf 4^4(0,3)$ tube (Section~\ref{rod1}):

\begin{center}
\begin{tabular}{l}
$S = \langle 3(00z), 2(x00), 2(x0\frac{1}{2}), t_x \rangle$;\\
$Stab_G(S) = \{e\}$;\\
$Stab_{N_{A}(G)}(S) = \langle 2(00z) \rangle$.
\end{tabular}
\end{center}

\noindent Here the computation of vertex stabilizer yields the same result as the conjugation action of the affine normalizer $N_A(P321) = P6/mmm \, ({\bf{a}}, {\bf{b}}, \frac{1}{2}{\bf{c}})$. Therefore, by adding the vertex stabilizer to $S$ we get: $Aut(\Gamma) \cong P622$, as expected from Theorem~\ref{teo1}, though the graph displays collisions in the barycentric placement.

\subsubsection{3-Periodic sphere-packing graphs}

Motivated by the above examples, we analyzed all three-periodic vertex-transitive sphere-packing graphs published by W.~Fischer and collaborators with respect to bounded automorphisms of finite order. To this end, using the computer algebra system $\GAP$, a code was written which -- given a generating set $S$ of a crystallographic group $G$ -- computes its stabilizer in $G$, $Stab_G(S)$, from a sufficiently large finite quotient $G/T^p \, (p \in \N)$ making use of the conjugacy separability property of space groups \cite{Dere26}. The code exploits the functionalities of $\GAP$ packages \emph{Cryst} \cite{Cryst2019} and \emph{Polycyclic}~\cite{Polyc2020}. The results are collected in Table~\ref{t:table4}. For the nomenclature of graphs see Appendix~B. The generating sets correspond to those listed by Fischer \cite{Fischer74, Fischer93} and Sowa \cite{Sowa2012}. $Stab_{Aut(\Gamma)}(x)$ was computed as described by Baburin (2020). The last column contains the isomorphism type of the quotient graph, $\Gamma/\sigma$, with respect to the partition $\sigma$ of $V(\Gamma)$ into orbits of $Stab_G(S)$. 


{\tiny
\begin{table}[H]
\begin{center}
\caption{Sphere-packing graphs with bounded automorphisms of finite order}
\label{t:table4}
\begin{tabular}{|l|l|c|c|c|}
\hline
Name & Gen. set $S$ & $Stab_G(S)$ & $Stab_{Aut(\Gamma)}(x)$ & $\Gamma/\sigma$ \\
\hline
\multicolumn{5}{|c|}{$I432$} \\
\hline
\rule{0pt}{2.5ex} $3/4/c3$  & $4(x00), 2(\frac{1}{2}-x, \frac{1}{4}, x)$ & $2(\frac{1}{2}y0)$ & $C_2 \times C_2$ & $3/4/c1$ \\
\rule{0pt}{2.5ex} $4/4/c13$ & $4(x00), 2(\frac{1}{2}+x, \frac{1}{4}, x), 2(\frac{1}{2}-x, \frac{1}{4}, x)$ & $2(\frac{1}{2}y0)$ & infinite & $''$ \\
\rule{0pt}{2.5ex} $4/4/c15$ & $4(x00), 2(\frac{1}{2}-x, \frac{1}{4}, x), 2(\frac{1}{2}y0)$ & $2(\frac{1}{2}y0)$ & $C_2 \times C_2$ & $''$ \\
\rule{0pt}{2.5ex} $5/3/c21$ & $4(x00), 2(\frac{1}{2}+x, \frac{1}{4}, x), 2(\frac{1}{2}y0),$ & $2(\frac{1}{2}y0)$ & infinite &  $''$ \\
& $2(\frac{1}{2}-x, \frac{1}{4}, x)$ & & & \\
\hline
\multicolumn{5}{|c|}{$I4_132$} \\
\hline
\rule{0pt}{2.5ex} $4/3/c22$          & $2(\frac{1}{4}y0), 3(xxx), 2(\frac{1}{4}+x, \frac{1}{8}, x)$ & $2(\frac{1}{4}-x, \frac{1}{8}, x)$ & infinite & $3/3/c1$ \\
\rule{0pt}{2.5ex} $\bolit {4/3/c25}$ & $3(xxx), 2(\frac{1}{4}+x, \frac{1}{8}, x), 2(\frac{1}{4}-x, \frac{1}{8}, x)$ & $2(\frac{1}{4}-x, \frac{1}{8}, x)$ & $C_2$ & $''$ \\
\rule{0pt}{2.5ex} $\bolit {4/3/c26}$ & $2(\frac{1}{4}y0), 3(xxx), 2(\frac{1}{4}-x, \frac{1}{8}, x)$ & $2(\frac{1}{4}-x, \frac{1}{8}, x)$ & $C_2$ & $''$ \\
\rule{0pt}{2.5ex} $5/3/c29$          & $2(\frac{1}{4}y0), 3(xxx), 2(\frac{1}{4}+x, \frac{1}{8}, x),$ & $2(\frac{1}{4}-x, \frac{1}{8}, x)$ & infinite & $''$ \\
\rule{0pt}{2.5ex} & $2(\frac{1}{4}-x, \frac{1}{8}, x)$ & & & \\
\hline
\multicolumn{5}{|c|}{$I4_122$} \\
\hline
\rule{0pt}{2.5ex} $\bolit {4/4/t43}$ & $4_3(\frac{1}{4}\frac{1}{4}z), 2(00z), 2(xx0)$ & $2(xx0)$ & $D_3$ & $3/10/c1$ \\
\rule{0pt}{2.5ex} $\bolit {5/4/t45}$ & $4_3(\frac{1}{4}\frac{1}{4}z), 4_1(\frac{1}{4}\frac{\bar1}{4}z), 2(\frac{1}{4}y\frac{3}{8})$ & $2(\frac{1}{4}y\frac{3}{8})$ & $C_2 \times C_2$ & $4/4/t1$ \\
\hline
\multicolumn{5}{|c|}{$I4_1/amd$} \\
\hline
\rule{0pt}{2.5ex} $\bolit {4/4/t53}$ & $4_3(\frac{1}{4}\frac{1}{4}z), m(x0z), 2(\frac{1}{4}y\frac{3}{8})$ & $2(\frac{1}{4}y\frac{3}{8})$ & $C_2$ & $3/8/t1$ \\
\hline
\multicolumn{5}{|c|}{$I\bar42d$} \\
\hline
\rule{0pt}{2.5ex} $\bolit {5/4/t45}$ & $d(x,\frac{1}{4}-x,z), d(x,x-\frac{1}{4},z), 2(\frac{1}{4}y\frac{3}{8})$ & $2(\frac{1}{4}y\frac{3}{8})$ & $C_2 \times C_2$ & $4/4/t1$ \\
\hline
\multicolumn{5}{|c|}{$I4_1/a$} \\
\hline
\rule{0pt}{2.5ex} $4/4/t39$ &  $\bar4_z(000), b(xy\frac{1}{8})$ & $2(00z)$ & infinite & $3/10/t4$ \\
\rule{0pt}{2.5ex} $5/3/t24$ &  $\bar4_z(000), b(xy\frac{1}{8}), 2(00z)$ & $2(00z)$ & infinite & $3/10/t4$ \\
\rule{0pt}{2.5ex} $\bolit {5/4/t45}$ &  $4_3(\frac{1}{4}\frac{1}{4}z), 4_1(\frac{1}{4}\frac{\bar1}{4}z), \bar1(\frac{1}{4}0\frac{3}{8})$ & $\bar1(\frac{1}{4}0\frac{3}{8})$ & $C_2 \times C_2$ & $4/4/t1$ \\
\hline
\multicolumn{5}{|c|}{$Fddd$} \\
\hline
\rule{0pt}{2.5ex} $\bolit {4/4/o18}$ &  $2(0y\frac{1}{2}), d(xy\frac{3}{8}), 2(00z)$ & $2(00z)$ & $C_2 \times C_2$ &  $3/10/t4$ \\
\rule{0pt}{2.5ex} $4/4/o19$ &  $2(\frac{1}{4}y\frac{1}{4}), 2(x\frac{1}{4}\frac{1}{4}), d(xy\frac{3}{8})$ & $2(\frac{1}{4}\frac{1}{4}z)$ & infinite & $3/10/t4$ \\
\rule{0pt}{2.5ex} $5/3/o6$  &  $2(\frac{1}{4}y\frac{1}{4}), 2(x\frac{1}{4}\frac{1}{4}), d(xy\frac{3}{8}), 2(\frac{1}{4}\frac{1}{4}z)$ & $2(\frac{1}{4}\frac{1}{4}z)$ & infinite & $3/10/t4$ \\
\hline
\end{tabular}
\end{center}
\end{table}
}

Out of 1080 sphere-packing graphs, there are 16 graphs\footnote{Note that graphs $\{4/3/c25, 4/3/c26\}$, $\{4/4/t39, 4/4/o19\}$, $\{5/3/t24, 5/3/o6\}$ are pairwise isomorphic.} which possess bounded automorphisms of finite order. Six of them are ladder graphs (highlighted in bold in Table~\ref{t:table4}), and others are related to them by addition or removal of some edges. It is noteworthy that almost half of the graphs have uncountable automorphism groups. Note that our results agree in part with those of Eon (2013) \cite{Eon2013} who considered $4/3/c25\text{-}26$, $4/4/o18\text{-}19$ and $5/3/o6/t24$ using the formalism of labelled quotient graphs.

In all the cases the stabilizer of the generating set has order~2, and therefore the orbits of (some subgroup of) $B_o(\Gamma)$ are always two-element blocks. Moreover, in most cases the stabilizer is contained in the generating set, and could be also determined from geometrical considerations. Qualitatively speaking, this stabilizer is an `echo' of centralizing elements in $Aut(\Gamma)$ which always correspond to bounded automorphisms \cite{Gromov83}.

Among the considered graphs, the sphere packing $3/8/t6$ (not in Table~\ref{t:table4}) deserves some comments. Fischer (1993) found it only in the general position of space group $I\bar42d$. The set of generators can be chosen as $\{2(x\frac{1}{4}\frac{1}{8}), d(x, \frac{1}{4}-x, z)\}$, the stabilizer of which is trivial. By Theorem~\ref{teo1}, $Aut(\Gamma)$ should be isomorphic to a space group. The graph has collisions in the barycentric placement that is, however, neighbour-unique. This allows us to get $Aut(\Gamma) \cong I4_1/amd$. The same result is obtained if the stabilizer of the generating set is referred to the affine normalizer of $I\bar42d$. The vertex stabilizer has order 2, and is unexpectedly (for valency 3) generated by the inversion through $(0\frac{1}{4}\frac{1}{8})$ forcing vertices to collide in $\R^3$ to the invariant position $I4_1/amd \, 8e \, (.2/m.)$. Therefore, regular groups of $3/8/t6$ must be noncentrosymmetric ($I\bar42d, I4_122, I4_1md, I\bar4m2$). Edge intersections are retained in the groups with mirror planes. Space groups $I\bar42d$ and $I4_122$ allow for piecewise-linear embeddings of the graph in $\R^3$. The first variant is realised with shortest equal distances and tabulated by Fischer (1993), the second [generators: $2(x\frac{1}{4}\frac{1}{8}), 2(xx0), 2(x,\frac{1}{2}+x, \frac{1}{4})$] does not correspond to a sphere packing. Both embeddings are compared in Fig.~\ref{f:beide}.

\begin{center}
\begin{figure}[h]
\centering
\includegraphics{Fig9.png}
%\includegraphics[width=18cm,height=8.5cm,keepaspectratio]{Fig9.png}
\caption{Embeddings of the $3/8/t6$ sphere-packing graph with different symmetry. The tetrahedral void at the center of the unit cell is emphasized in both structures. Notice also different point-group symmetry of the tetrahedra ($\bar4$ \emph{versus} $222$).}
\label{f:beide}
\end{figure}
\end{center}
\vspace{-2em}

The pair of sphere packings $4/3/c25\text{-}26$ is remarkable as they arise from non-isotopic embeddings of one and the same ladder graph built from the two copies of $3/3/c1$. This raises a natural question -- given a three-periodic vertex-transitive ladder graph, how to find its (piecewise linear) vertex-transitive embeddings in $\R^3$? A practical algorithm for this problem is suggested in the next section.


\subsubsection{Construction of 3-periodic \emph{ladder graphs} from $\R^4$}

A suitable (piecewise-linear) embedding of a periodic graph in $\R^d$ may be obtained in two steps: (a) compute a symmetry group of the embedding based on the knowledge of $Aut(\Gamma)$, (b) given a symmetry group, find coordinates which are compatible with it and are optimal in some sense. The first step is group-theoretic, the second is numerical. It is the first step that we address here for 3-periodic vertex-transitive ladder graphs in $\R^3$ (this topic has been considered recently by O'Keeffe \emph{et al.} \cite{OK25} in an essentially empirical way). The second step is carried out by setting up a system of quadratic equations for the distances to be solved usually by least-squares methods.

We start with a 3-periodic vertex-transitive graph $\Gamma$ with nontrivial vertex-stabilizers in $Aut(\Gamma)$ that is isomorphic to a space group in $\R^3$, lift it to a 3-dimensional hyperplane of $R^4$, say, $w = 1/5$, and the reflection in the hyperplane $w = 0$ generates a second copy of it. In this way a 3-periodic ladder graph $\Lambda$ is obtained whose $Aut(\Lambda) \cong Aut(\Gamma) \times C_2$, and vertex stabilizers are isomorphic to those of the parent graph $\Gamma$ in $Aut(\Gamma)$. Now, compute proper vertex-transitive subgroups of $Aut(\Lambda)$ as subperiodic groups in $\R^4$ (with a three-dimensional translation lattice), and project them back to $\R^3$ by restricting symmetry operators to the respective subspace. The groups not containing the reflection at $w=0$ are precisely space groups for vertex-transitive embeddings of $\Lambda$ in $\R^3$.

\medskip

\noindent \emph{Example}. Construct a ladder graph for $\Gamma = 3/3/c1$ using the proposed approach. We have:

\begin{center}
\begin{tabular}{l}
$Aut(\Gamma) = I4_132, Stab_{Aut(\Gamma)}(x) \cong C_2$;\\
$Aut(\Lambda) \cong I4_132 \times C_2, Stab_{Aut(\Lambda)}(x) \cong C_2$.
\end{tabular}
\end{center}

\noindent It can be shown by computation that the group $Aut(\Lambda)$ has two subgroups of index~2, $G_1$ and $G_2$, without a mirror plane at $w=0$ (\emph{cf.}~above). They are the two vertex-transitive (and in this case necessarily regular) subgroups of $Aut(\Lambda)$. In terms of group structure, $Aut(\Lambda)$, $G_1$ and $G_2$ are quotients of four-dimensional space groups $25/10/2/3$, $24/4/3/3$ and $25/10/1/8$ (in the notation of \cite{BBNWZ}), respectively, obtained by factoring out the normal translation subgroup of rank~1. The point groups are $P(G_1) \cong S_4$ and $P(Aut(\Lambda)) = P(G_2) \cong S_4 \times C_2$, hence $G_1$ and $G_2$ are \emph{translationengleiche} and \emph{klassengleiche} subgroups of $Aut(\Lambda)$, resp. Upon restricting symmetry operators to $\R^3$, the glide reflection of $G_2$, $\{m_{xyz}|\frac{1}{2}\frac{1}{2}\frac{1}{2}0\}$, is mapped to the $I$-translation of $I4_132$ with a simultaneous reduction of the point group, and finally we get $G_1 \cong G_2 \cong I4_132$. Note that $G_1$ corresponds to $4/3/c26$ and $G_2$ to $4/3/c25$. They are illustrated in Fig.~\ref{f:ulds}.

%\medskip

\begin{center}
\begin{figure}[h]
\centering
\includegraphics{Fig10.png}
\caption{$4/3/c25$ and $4/3/c26$ sphere-packing graphs. Thin black edges interlink the two connected components.}
\label{f:ulds}
\end{figure}
\end{center}
\vspace{-2em}


To demonstrate the capability of the proposed approach, in Table~\ref{t:table5} regular automorphism groups $G(\Lambda)$ are given for ladder graphs built from selected graphs of different symmetry (only regular groups are listed in order to save space, and moreover, they are helpful for understanding the method). Therein, the last element of a generating set in the fourth column is always one exchanging the two 3-periodic subgraphs isomorphic to $\Gamma$, \emph{i.e.} it is the element of the coset $G(\Lambda)/Stab_{G(\Lambda)}(\Gamma)$. Observe that this coset must not contain mirror reflections because otherwise edge intersections are enforced in the embedding \cite{Bab16}. For example, with symmetry $Pn\bar3m$ the ladder graph can be constructed from $3/4/c1$ in two different ways. Consider the variant corresponding to the first entry of Table~\ref{t:table5}. In this case $Stab_G(\Gamma) \cong Fd\bar3c$, \emph{i.e.} the group without diagonal mirrors of $G=Pn\bar3m$. As a result, this generating set always leads to edge intersections in the embedding of $\Lambda$.

{\tiny
\begin{table}[H]
\begin{center}
\caption{Examples of ladder-graph construction}
%\vspace{-0.8em}
\label{t:table5}
\begin{tabular}{|c|l|l|l|l|}
\hline
$\Gamma$ & $G(\Lambda)$ & $Stab_{G(\Lambda)}(\Gamma)$ & Generators for $G(\Lambda)$ & Source/ \\
& & & & related graph \\
\hline
\rule{0pt}{2.5ex} $3/4/c1$ & $Pn\bar3m$ & $Fd\bar3c \, (2{\bf{a}})$ & $\bar4_z(00\frac{1}{2}), 2(\frac{1}{4},y,\frac{1}{2}+y), 2(x0\frac{1}{2})$ & intersect. \\
	 &            & $Fd\bar3m \, (2{\bf{a}})$ & $m(xxz), m(x\bar{x}z), 2(\frac{1}{4}, y, \frac{\bar1}{2}-y), 2(x0\frac{\bar1}{2})$ & \cite{Blatov07} \\
$3/8/h1$ & $R\bar3m$  & $R3m$ & $3_2(\frac{1}{3}0z), m(x, 2x-1, z), 2(x00)$ & -- \\
	 &            & $R\bar3m \, (2{\bf{c}})$ & $2(x,x-\frac{1}{3},\frac{1}{3}), 2(\frac{1}{3}y\frac{2}{3}), m(x, 2x-1, z), 2(x0\frac{1}{2})$ & \cite{Blatov07} \\
$3/8/t1$ & $I4_1/amd$ & $I4_1md$ & $4_3(\frac{1}{4}\frac{1}{4}z), m(0yz), 2(x\frac{1}{4}\frac{1}{8})$ & $4/4/t53$ \\
         &            & $I\bar4m2$ & $2(xx0), m(0yz), 2(x, \frac{1}{2}-x, \frac{1}{4}), 2(x\frac{1}{4}\frac{1}{8})$ & -- \\
\rule{0pt}{2.5ex} $3/10/c1$ & $I4_122$  & $I2_12_12_1$ & $2(x\frac{1}{4}\frac{\bar3}{8}), 2(\frac{1}{4}y\frac{\bar1}{8}), 2(\frac{1}{2}0z), 2(x, \frac{1}{2}-x, \frac{\bar1}{4})$ & $4/3/c26$ \\
\rule{0pt}{2.5ex}  &  & $I4_1$ & $4_3(\frac{1}{4}\frac{1}{4}z), 2(\frac{1}{2}0z), 2(x, \frac{1}{2}-x, \frac{\bar1}{4})$ & $4/4/t43$ \\
          &           & $P4_32_12$ & $4_3(\frac{1}{4}\frac{1}{4}z), 2(x, x-\frac{1}{2}, \frac{\bar1}{4}), 2(x, \frac{1}{2}-x, \frac{\bar1}{4})$ & $4/3/c25$ \\
\rule{0pt}{2.5ex} &   & $P4_122$ & $2(x\frac{1}{4}\frac{\bar3}{8}), 2(\frac{1}{4}y\frac{\bar1}{8}), 2(x, x-\frac{1}{2}, \frac{\bar1}{4}), 2(x, \frac{1}{2}-x, \frac{\bar1}{4})$ & -- \\
\rule{0pt}{2.5ex} $3/12/h1$ & $P6_222$  & $P6_2$   & $6_2(00z), 2(\frac{1}{2}0z), 2(x00)$ & -- \\
          &           & $P3_212$ & $2(x\bar{x}\frac{\bar1}{6}), 2(2x,x,\frac{1}{6}), 2(x, 2x-1, 0), 2(x00)$ & -- \\
          &           & $P6_422 \, (2{\bf{c}})$ & $2(\frac{1}{2}0z), 2(x\bar{x}\frac{1}{3}), 2(2x, x, \frac{2}{3}), 2(x0\frac{1}{2})$ & -- \\
          &           & $P6_122 \, (2{\bf{c}})$ & $6_2(00z), 2(x, 2x-1, \frac{1}{2}), 2(x0\frac{1}{2})$ & -- \\
\hline
\end{tabular}
\end{center}
\end{table}
}

Table~\ref{t:table5} shows that known ladder graphs are reproduced, and more are found. For example, a new embedding is derived for the ladder graph of $3/10/c1$ (see coordinates in the Supporting information). The ladder graphs for $3/4/c1$ and $3/8/h1$ from Blatov (2007) \cite{Blatov07} are reproduced which (in his nomenclature) should be named {\bf nbo-g}-4-$Pn\bar3m$ and {\bf smd}-4-$R\bar3m$, resp. Our method also allows to clarify why some graphs could not be encountered in former studies: for example, the ladder graph of $3/6/h1$ (a cubic lattice with nodes replaced by hexagons) admits only embeddings with edge intersections because mirror reflections always belong to the cosets $G(\Lambda)/Stab_{G(\Lambda)}(\Gamma)$.

\section{Appendix A: Some more group theory}

There exists a deep connection between the structure of a graph and the subgroup structure of its automorphism group. The asymptotic behaviour of the growth function of a graph (coordination sequence in crystallography) is the property that determines this connection. For vertex-transitive graphs, the problem can be put into a group-theoretic setting as described below (here it is enough to keep in mind that the growth of a group is equivalent to that of its Cayley graph).

\begin{theorem}
\label{teo3}
\emph{(Gromov, 1981)}. If $G$ is a finitely generated group with polynomial growth, then it contains a nilpotent subgroup of finite index.
\end{theorem}

\noindent This means that if $\Gamma$ is a Cayley graph of $G$, then a nilpotent subgroup of $G$ acts on $V(\Gamma)$ with a finite number of orbits. Note that the class of torsion-free nilpotent groups of growth rate $d$ coincide with that of free abelian groups $\Z^d$ for $d \leq 3$ \cite{Seif97}. However, not all graphs are Cayley graphs, \emph{i.e.} not all allow for a regular group of automorphisms. The following theorem covers this case and, more importantly, is applicable to graphs with uncountable automorphism groups.

\begin{theorem}
\label{teo4}
\emph{(Trofimov, 1984)}. Let $\Gamma$ be a vertex-transitive graph with polynomial growth. Then there exists a partition $\sigma$ of $V(\Gamma)$ into finite blocks $x^\sigma$ such that the group $Aut(\Gamma/\sigma)$ is finitely generated and contains a nilpotent subgroup of finite index. Moreover, $Stab_{Aut(\Gamma/\sigma)}(x^\sigma)$ is finite.
\end{theorem}

\noindent Theorem~\ref{teo5} uncovers the only `mechanism' for graphs of polynomial growth to have uncountable $Aut(\Gamma)$: it is the possibility to permute vertices inside blocks $x^\sigma$ independently (such as the blocks $\{u_i, v_i\}$ in Fig.~\ref{f:infstab}).

Let $\Gamma$ be a Cayley graph of a space group $G$ and let $Aut(\Gamma)$ be such that $B_o(\Gamma) \neq \{e\}$ (for the notation see Section~4.2). Based on \cite{Trofimov83, Trofimov84}, we have the following general result.

\begin{theorem}
\label{teo5}
Under the assumptions made, $B_o(\Gamma) \in C_{Aut(\Gamma)}(T(G))$ and hence $B_o(\Gamma) \in C_{Aut(\Gamma)}(G)$.
\end{theorem}

\noindent If $|Aut(\Gamma):G|$ is finite (and equal to $|Stab_{Aut(\Gamma)}(x)|$), then $Aut(\Gamma)$ is virtually abelian and $B(\Gamma)$ is the FC-center of $Aut(\Gamma)$.

\section{Appendix B: Nomenclature for periodic graphs}
For the nomenclature of periodic graphs, the notation proposed in \cite{Fischer71, KF78} is preferred for many reasons (one of them is that most subperiodic graphs are not included, for example, in the popular RCSR database \cite{RCSR}). 1-Periodic graphs are characterized by the Schl{\"a}fli symbol of a vertex-transitive planar graph followed by the rolling vector (referred to the usual choice of basis vectors \cite{ITE}). For example, the famous Boerdijk--Coxeter helix would be called $\bf 3^6(1, 3)$. For three-periodic sphere-packing graphs, the symbols used are of type $k/m/fn$: here $k$ stands for the valency, $m$ for the girth (shortest cycle) of a graph, $f$ refers to the crystal system (where the corresponding sphere packing is realized with highest symmetry), and $n$ is an arbitrary number. Table~\ref{t:table6} contains correspondences (if available) between original Fischer names and the RCSR three-letter codes. %The latter can be used to search for the graph data in the online database  \cite{RCSR}.


{\tiny
\begin{table}[H]
\begin{center}
\caption{Fischer vs. RCSR}
\label{t:table6}
\begin{tabular}{|c|c|}
\hline
Fischer & RCSR  \\
\hline
$3/3/c1$  & {\bf srs-a} \\
$3/4/c1$  & {\bf nbo-a} \\
$3/10/c1$ & {\bf srs} \\
$4/3/c25$ & {\bf uld} \\
$4/3/c26$ & {\bf uld-z} \\
$4/4/c15$ & {\bf ulm} \\
$5/3/c21$ & {\bf nbo-g} \\
$5/3/c29$ & {\bf fco} \\
$4/6/c1$  & {\bf dia} \\
$6/4/c1$  & {\bf pcu} \\
$3/8/t1$  & {\bf lig} \\
$3/10/t4$ & {\bf ths} \\
$4/4/t1$  & {\bf lvt} \\
$4/6/t4$  & {\bf cds} \\
$3/6/h1$ & {\bf pcu-h} \\
$3/8/h1$  & {\bf etb} \\
$3/12/h1$ & {\bf twt} \\
\hline
\end{tabular}
\end{center}
\end{table}
}


%\bigskip

\noindent {\bf Acknowledgement}. Part of this work was done in 2022--2023 when the author served as a substitute professor (Vertretungsprofessor) at the Ludwig-Maximilians-Universit\"at M\"unchen, Sektion Kristallographie. Funding within the Hightech Agenda Bayern is gratefully acknowledged.

\begin{thebibliography}{99}

\setlength\itemsep{-0.01cm}

\begin{small}

\bibitem{Bab16} Baburin, I.~A. On the group-theoretical approach to the study of interpenetrating nets. Acta Cryst. 2016, A72, 366-375.

\bibitem{Bab20} Baburin, I.~A. On Cayley graphs of $\Z^4$. Acta Cryst. 2020, A76, 584-588.

\bibitem{Bab26} Baburin, I.~A. Short presentations for crystallographic groups. Acta Cryst. 2026, A82, 18-31.

\bibitem{Evar09} Bandura, A.~V. and Evarestov, R.~A. From anatase $(101)$ surface to TiO\textsubscript{2} nanotubes: rolling procedure and first-principles LCAO calculations. Surface Science, 2009, 603, L117-L120.

\bibitem{Besche2002} Besche, H.~U., Eick B., O'Brien, E.~A. A millenium project: constructing small groups. Int. J. Algebra Comput., 2002, 12, 623-644.

\bibitem{Biggs93} Biggs, N. Algebraic Graph Theory, 2nd edition. Cambridge University Press, 1993.

\bibitem{Blatov07} Blatov V.~A. Topological relations between three-dimensional periodic nets. I. Uninodal nets. Acta Cryst. 2007, A63, 329-343.

\bibitem{BBNWZ} Brown, H., B{\"u}low, R., Neub{\"u}ser, J., Wondratschek, H., Zassenhaus, H. Crystallographic groups of four-dimensional space. Wiley, 1978.

\bibitem{Proser03} Carlucci, L., Ciani, G. and Proserpio, D.~M. Polycatenation, polythreading and polyknotting in coordination network chemistry. Coord. Chem. Rev., 2003, 246, 247-289.

\bibitem{Linegr10}  Damnjanovi{\'c}, M., Milo{\v{s}}sevi{\'c}, I. Line groups in physics. Springer, 2010.

\bibitem{Delgado2004} Delgado-Friedrichs, O. Barycentric drawings of periodic graphs. Lecture Notes in Computer Science, vol. 2912. Springer, 2004, 178-189.

\bibitem{OK25} Delgado-Friedrichs, O., O'Keeffe, M. and Treacy, M.~M.~J. Periodic graphs with coincident edges: folding-ladder and related graphs. Acta Cryst. 2025, A81, 49-56.

\bibitem{Dere26} Der\'e, J. and Vandeputte, L. Effective conjugacy separability of virtually abelian groups. Comm. in Algebra, 2026, 54, 1129-1152.

\bibitem{Cryst2019} Eick, B., G{\"a}hler, F., Nickel, W. \emph{Cryst} -- computing with crystallographic groups, a refereed $\GAP$~4 package.

\bibitem{Polyc2020} Eick, B., Nickel, W., Horn, M. \emph{Polycyclic} -- computation with polycyclic groups, a refereed $\GAP$~4 package.

\bibitem{Eon2005} Eon, J.-G. Graph-theoretical characterization of periodicity in crystallographic nets and other infinite graphs. Acta Cryst. 2005, A61, 501-511.

\bibitem{Fischer71} Fischer, W. Existenzbedingungen homogener Kugelpackungen in Raumgruppen tetragonaler Symmetrie. Z. Kristallogr., 1971, 133, 18-42.

\bibitem{Fischer74} Fischer, W. Existenzbedingungen homogener Kugelpackungen zu kubischen Gitterkomplexen mit drei Freiheitsgraden. Z.~Kristallogr., 1974, 140, 50-74.

\bibitem{Fischer93} Fischer, W. Tetragonal sphere packings. III. Lattice complexes with three degrees of freedom. Z.~Kristallogr., 1993, 205, 9-26.

\bibitem{GAP} The $\GAP$ Group. $\GAP$ -- Groups, Algorithms, Programming. Version 4.19. Available from https://www.gap-system.org.

\bibitem{Gods80} Godsil, C.~D. and McKay, B.~D. The dimension of a graph. Quart. J. Math. Oxford, 1980, 31, 423-427.

\bibitem{Seif89all} Godsil, C.~D., Imrich, W., Seifter, N., Watkins, M.~E. and Woess, W. A note on bounded automorphisms of infinite graphs.  Graphs and Combinatorics, 1989, 5, 333-338.

\bibitem{Gromov81} Gromov, M. Groups of polynomial growth and expanding maps. Publ. Math. Inst. Hautes \'{E}tudes Sci., 1981, 53, 53-78.

\bibitem{Gromov83} Gromov, M. Infinite groups as geometric objects. ICM Proceedings, Warszawa, 1983, Vol.~1, 385-392.

%\bibitem{Seif88} Imrich, W. and Seifter, N. A note on the growth of transitive graphs. Discrete Math., 1988/89, 73, 111-117.

\bibitem{Seif91} Imrich, W. and Seifter, N. A survey on graphs with polynomial growth. Discrete Math., 1991, 95, 101-117.

\bibitem{ITE} International tables for crystallography. Volume E (2006). Available from https://it.iucr.org.

\bibitem{Jones11} Jones, G.~A. Detecting and identifying wallpaper groups. Symmetry: Culture and Science, 2011, 22, pp. 131-157.

\bibitem{Jung84} Jung, H.~A. and Watkins, M.~E. Fragments and automorphisms of infinite graphs. Eur. J. Comb., 1984, 5, 149-162.

\bibitem{Twist2012} Kato, K., Koretsune, T. and Saito, S. Energetics and electronic properties of twisted single-walled carbon nanotubes. Phys.~Rev.~B 85, 115448.

\bibitem{KF75} Koch, E. and Fischer, W. Automorphismengruppen von Raumgruppen und die Zuordnung von Punktlagen zu Konfigurationslagen. Acta Cryst. 1975, A31, 88-95.

\bibitem{KF78b} Koch, E. and Fischer, W. Complexes for crystallographic point groups, rod groups and layer groups. Z. Kristallogr., 1978, 147, 21-38.

\bibitem{KF78} Koch, E. and Fischer, W. Types of sphere packings for crystallographic point groups, rod groups and layer groups. Z. Kristallogr., 1978, 148, 107-152.

\bibitem{Koe80} K{\"o}hler, K.-J. On the structure and the determination of $n$-dimensional partially periodic crystallographic groups. \emph{MATCH}, 1981, 10, 27-53.

\bibitem{layer25} Mahmoudi, S., Dresselhaus, E.~J. and Dimitriyev, M.~S. An orbifold framework for classifying layer groups with an application to knitted fabrics. arXiv:2512.05149.

\bibitem{Eon2013} Moreira de Oliveira Jr, M. and Eon, J.-G. Non-crystallographic nets with finite blocks of imprimitivity for bounded automorphisms.  Acta Cryst. 2013, A69, 276-288.

\bibitem{RCSR} O'Keeffe, M., Peskov, M.~A., Ramsden, S.~J. and Yaghi, O.~M. The Reticular Chemistry Structure Resource (RCSR) database of, and symbols for, crystal nets. Acc. Chem. Res. 2008, 41, 1782-1789. http://rcsr.net

\bibitem{OK21} O'Keeffe, M. and Treacy, M.~M.~J. Isogonal piecewise linear embeddings of 1-periodic weaves and some related structures. Acta Cryst. 2021, A77, 130-137.

\bibitem{Salle19} de la Salle, M., Tessera, R. Characterizing a vertex-transitive graph by a large ball. Journal of topology, 2019, 12, 704-742.

\bibitem{Tes21} Tessera, R. and Tointon, M.~C.~H. A finitary structure theorem for vertex-transitive graphs of polynomial growth. Combinatorica, 2021, 41, 263-298.

\bibitem{Seif97} Seifter, N. and Trofimov, V.~I. Automorphism groups of graphs with quadratic growth. J. Comb. Theory B, 1997, 71, 205-210.

\bibitem{Souv26} Souvignier, B. It's all in the group. Acta Cryst., 2026, A82, 79-82. 

\bibitem{Sowa2012} Sowa, H. Orthorhombic sphere packings. IV. Trivariant lattice complexes of space groups without mirror  planes belonging to crystal class $mmm$.  Acta Cryst., 2012, A68, 763-777.

\bibitem{Tom84} Tomkinson, M.~J. FC-groups, Res. Notes in Math. 96. Pitman, London, 1984.

\bibitem{Trofimov83} Trofimov, V.~I. Automorphisms of graphs and a characterization of lattices. Math. USSR Izvestiya, 1984, 22, 379-391. % https://www.mathnet.ru/eng/im1408

\bibitem{Trofimov84} Trofimov, V.~I. Graphs with polynomial growth. Math. USSR Sbornik, 1985, 51, 405-417. % https://www.mathnet.ru/eng/sm2028

\bibitem{Watk71} Watkins, M.~E. On the action of non-abelian groups on graphs. J. Comb. Theory, 1971, 11, 95-104.

\bibitem{Wie64} Wielandt, H. Finite permutation groups. Academic Press, 1964.

\end{small}

\end{thebibliography}


\end{document}
